\documentclass[11pt]{article}

\usepackage[a4paper,margin=1in]{geometry}
\usepackage{amsmath,amssymb,amsthm,mathtools}
\usepackage[T1]{fontenc}
\usepackage{lmodern}
\usepackage{microtype}

\newtheorem{theorem}{Theorem}[section]
\newtheorem{lemma}[theorem]{Lemma}
\newtheorem{proposition}[theorem]{Proposition}
\newtheorem{remark}[theorem]{Remark}

\newcommand{\C}{\mathbb{C}}
\newcommand{\R}{\mathbb{R}}
\newcommand{\eps}{\varepsilon}

\title{The Real-Slice Criterion for the Clean OBC Hatano--Nelson Family\\
in Dimensions \(2\) and \(3\)}
\author{}
\date{}

\begin{document}
\maketitle

\section{Setup}

Fix real numbers \(0<a<b\), and let
\[
\mathbf A_n=\operatorname{tridiag}(a,0,b)
\in \R^{n\times n}.
\]
We are interested here only in the cases \(n=2\) and \(n=3\):
\[
\mathbf A_2=\begin{pmatrix}0&b\\ a&0\end{pmatrix},
\qquad
\mathbf A_3=\begin{pmatrix}0&b&0\\ a&0&b\\ 0&a&0\end{pmatrix}.
\]

For a matrix \(\mathbf A\in\C^{N\times N}\) and \(\eps>0\), its \(\eps\)-pseudospectrum is
\[
\sigma_\eps(\mathbf A)
\coloneqq
\bigl\{z\in\C:\|(z\mathbf I-\mathbf A)^{-1}\|>\eps^{-1}\bigr\},
\]
where \(\|\cdot\|\) is the operator norm induced by the Euclidean norm, with the convention
\[
\|(z\mathbf I-\mathbf A)^{-1}\|=+\infty
\qquad\text{for } z\in\sigma(\mathbf A).
\]

For \(x,y\in\R\), define
\[
\mathbf M_n(x,y)
\coloneqq
\bigl((x+iy)\mathbf I-\mathbf A_n\bigr)^*
\bigl((x+iy)\mathbf I-\mathbf A_n\bigr),
\]
and let
\[
\lambda_n(x,y)\coloneqq \lambda_{\min}\!\bigl(\mathbf M_n(x,y)\bigr)
=
\sigma_{\min}\!\bigl((x+iy)\mathbf I-\mathbf A_n\bigr)^2.
\]
Therefore
\[
x+iy\in \sigma_\eps(\mathbf A_n)
\quad\Longleftrightarrow\quad
\lambda_n(x,y)<\eps^2.
\]

Throughout, we write
\[
p\coloneqq a+b,
\qquad
q\coloneqq b-a>0.
\]

\section{A topological reduction}

The following elementary lemma isolates the geometric content that we need.

\begin{lemma}\label{lem:topological}
Let \(\Omega\subset\C\) be open. Assume that for every \(x\in\R\), the vertical section
\[
V_x(\Omega)\coloneqq \{y\in\R:\ x+iy\in\Omega\}
\]
is either empty or an open interval containing \(0\). Then the connected components of
\(\Omega\) are in one-to-one correspondence with the connected components of
\(\Omega\cap\R\). In particular,
\[
\Omega\ \text{is connected}
\quad\Longleftrightarrow\quad
\Omega\cap\R\ \text{is connected}.
\]
\end{lemma}

\begin{proof}
Set
\[
\Sigma\coloneqq \Omega\cap\R.
\]
Let \(I\) be a connected component of \(\Sigma\), and define
\[
\Omega_I\coloneqq \{z\in\Omega:\ \operatorname{Re} z\in I\}.
\]
We claim that \(\Omega_I\) is path connected.

Indeed, let
\[
z_1=x_1+iy_1\in \Omega_I,
\qquad
z_2=x_2+iy_2\in \Omega_I.
\]
Since \(x_j\in I\) and the vertical section at \(x_j\) contains \(0\), the vertical segment
joining \(z_j\) to \(x_j\) lies in \(\Omega\). Also, because \(I\subset \Sigma\subset \Omega\),
the real line segment joining \(x_1\) to \(x_2\) lies in \(\Omega\). Hence \(z_1\) and \(z_2\)
can be joined by a polygonal path contained in \(\Omega_I\).

Now let \(I_1\neq I_2\) be distinct connected components of \(\Sigma\). Then they are separated by
a nonempty open interval \(J\subset\R\) with \(J\cap \Sigma=\varnothing\). We claim that the
vertical strip
\[
S_J\coloneqq \{z\in\C:\ \operatorname{Re} z\in J\}
\]
is disjoint from \(\Omega\). If not, some \(x+iy\in \Omega\) would satisfy \(x\in J\). Then
\(V_x(\Omega)\) would be nonempty, hence would contain \(0\), so \(x\in \Sigma\), contradicting
\(x\in J\subset \R\setminus\Sigma\). Thus \(S_J\cap\Omega=\varnothing\), so \(\Omega_{I_1}\) and
\(\Omega_{I_2}\) lie in distinct connected components of \(\Omega\).

Therefore the sets \(\Omega_I\), as \(I\) ranges over the connected components of \(\Sigma\), are
exactly the connected components of \(\Omega\). This proves the lemma.
\end{proof}

\section{The case \(n=2\)}

\begin{proposition}\label{prop:n2}
Fix \(x\in\R\). Then \(\lambda_2(x,y)\) depends only on \(y^2\), and is a nondecreasing function
of \(y^2\) on \([0,\infty)\). Equivalently,
\[
|y|\longmapsto \sigma_{\min}\!\bigl((x+iy)\mathbf I-\mathbf A_2\bigr)
\]
is nondecreasing on \([0,\infty)\).
\end{proposition}

\begin{proof}
Set
\[
r\coloneqq x^2+y^2.
\]
A direct computation gives
\[
\mathbf M_2(x,y)
=
\begin{pmatrix}
r+a^2 & -px+iqy\\
-px-iqy & r+b^2
\end{pmatrix}.
\]
Hence
\[
\lambda_2(x,y)
=
r+\frac{a^2+b^2}{2}
-\frac12\sqrt{(a^2-b^2)^2+4p^2x^2+4q^2y^2}.
\]
Since
\[
a^2-b^2=-(a+b)(b-a)=-pq,
\]
this may be rewritten as
\[
\lambda_2(x,y)
=
x^2+y^2+\frac{a^2+b^2}{2}
-\frac12\sqrt{p^2q^2+4p^2x^2+4q^2y^2}.
\]
Thus \(\lambda_2(x,y)\) depends only on \(y^2\). Writing \(t=y^2\), we get
\[
\lambda_2(x,\sqrt t)
=
x^2+t+\frac{a^2+b^2}{2}
-\frac12\sqrt{p^2q^2+4p^2x^2+4q^2t},
\]
and therefore
\[
\frac{d}{dt}\lambda_2(x,\sqrt t)
=
1-\frac{q^2}{\sqrt{p^2q^2+4p^2x^2+4q^2t}}.
\]
Since
\[
\sqrt{p^2q^2+4p^2x^2+4q^2t}\ge pq>q^2
\]
(because \(p=a+b>b-a=q\)), it follows that
\[
\frac{d}{dt}\lambda_2(x,\sqrt t)\ge 0.
\]
Hence \(\lambda_2(x,y)\) is a nondecreasing function of \(y^2\), as claimed.
\end{proof}

\section{The case \(n=3\)}

\begin{proposition}\label{prop:n3}
Fix \(x\in\R\). Then \(\lambda_3(x,y)\) depends only on \(y^2\), and is a nondecreasing function
of \(y^2\) on \([0,\infty)\). Equivalently,
\[
|y|\longmapsto \sigma_{\min}\!\bigl((x+iy)\mathbf I-\mathbf A_3\bigr)
\]
is nondecreasing on \([0,\infty)\).
\end{proposition}

\begin{proof}
Set
\[
r\coloneqq x^2+y^2.
\]
Introduce the orthonormal basis
\[
e_+\coloneqq \frac1{\sqrt2}(1,0,1)^T,
\qquad
e_0\coloneqq (0,1,0)^T,
\qquad
e_-\coloneqq \frac1{\sqrt2}(1,0,-1)^T,
\]
and let \(\mathbf P\) be the orthogonal matrix with columns \(e_+,e_0,e_-\). A direct computation gives
\[
\mathbf P^* \mathbf M_3(x,y)\mathbf P=
\begin{pmatrix}
r+\dfrac{p^2}{2} & -\sqrt2\,p\,x & -\dfrac{pq}{2}\\[1ex]
-\sqrt2\,p\,x & r+\dfrac{p^2+q^2}{2} & -i\sqrt2\,q\,y\\[1ex]
-\dfrac{pq}{2} & i\sqrt2\,q\,y & r+\dfrac{q^2}{2}
\end{pmatrix}.
\]
Thus \(\lambda_3(x,y)\) is the smallest eigenvalue of this Hermitian matrix.

Since
\[
\mathbf A_3\,(b,0,-a)^T=0,
\]
we have
\[
\lambda_3(x,y)\le
\frac{\|(x+iy)(b,0,-a)^T\|^2}{\|(b,0,-a)^T\|^2}
=
x^2+y^2=r.
\]
Hence
\[
\nu\coloneqq r-\lambda_3(x,y)\ge 0.
\]
Define
\[
\alpha\coloneqq \nu+\frac{p^2}{2},
\qquad
\beta\coloneqq \nu+\frac{p^2+q^2}{2},
\qquad
\gamma\coloneqq \nu+\frac{q^2}{2}.
\]
Then
\[
\mathbf P^* \mathbf M_3(x,y)\mathbf P-\lambda_3(x,y)\mathbf I
=
\begin{pmatrix}
\alpha & -\sqrt2\,p\,x & -\dfrac{pq}{2}\\[1ex]
-\sqrt2\,p\,x & \beta & -i\sqrt2\,q\,y\\[1ex]
-\dfrac{pq}{2} & i\sqrt2\,q\,y & \gamma
\end{pmatrix}.
\]
Since \(\lambda_3(x,y)\) is an eigenvalue, the determinant of this matrix is zero. Expanding
the determinant gives
\[
\alpha(\beta\gamma-2q^2y^2)-2p^2x^2\gamma-\frac{p^2q^2}{4}\beta=0.
\]
Now
\[
\alpha\gamma-\frac{p^2q^2}{4}
=
\left(\nu+\frac{p^2}{2}\right)\left(\nu+\frac{q^2}{2}\right)-\frac{p^2q^2}{4}
=
\nu\left(\nu+\frac{p^2+q^2}{2}\right)
=
\nu\beta,
\]
so the determinant equation simplifies to
\begin{equation}\label{eq:n3-master}
\nu\,\beta^2=2q^2y^2\,\alpha+2p^2x^2\,\gamma.
\end{equation}

Now set
\[
t\coloneqq y^2.
\]
Because \(\lambda_3(x,y)\) depends continuously on \(y\), the function
\[
\nu(t)\coloneqq x^2+t-\lambda_3(x,\sqrt t)
\]
is continuous on \([0,\infty)\), and its image
\[
I\coloneqq \nu([0,\infty))
\]
is an interval in \([0,\infty)\).

For \(\nu\in I\), equation \eqref{eq:n3-master} determines \(t\) as
\[
t=\Phi(\nu)\coloneqq
\frac{\nu\,\beta(\nu)^2-2p^2x^2\,\gamma(\nu)}{2q^2\,\alpha(\nu)},
\]
where
\[
\alpha(\nu)=\nu+\frac{p^2}{2},
\qquad
\beta(\nu)=\nu+\frac{p^2+q^2}{2},
\qquad
\gamma(\nu)=\nu+\frac{q^2}{2}.
\]

We claim that
\[
\Phi'(\nu)\ge 1
\qquad\text{for all } \nu\in I.
\]
Indeed, differentiating \(\Phi\) gives
\[
2q^2\alpha\,\Phi'
=
\beta^2+2\nu\beta-2q^2\Phi-2p^2x^2.
\]
Since \(\Phi=t\) on \(I\), and \eqref{eq:n3-master} says
\[
2q^2t=\frac{\nu\beta^2-2p^2x^2\gamma}{\alpha},
\]
we obtain
\[
2q^2\alpha\,\Phi'
=
\beta^2+2\nu\beta-\frac{\nu\beta^2}{\alpha}
-\frac{p^2x^2(p^2-q^2)}{\alpha}.
\]
Also, \eqref{eq:n3-master} implies
\[
2p^2x^2\gamma\le \nu\beta^2,
\]
hence
\[
p^2x^2\le \frac{\nu\beta^2}{2\gamma}.
\]
Therefore
\[
2q^2\alpha\,\Phi'
\ge
\beta^2+2\nu\beta-\frac{\nu\beta^2}{\alpha}
-\frac{\nu\beta^2(p^2-q^2)}{2\alpha\gamma}.
\]
Since
\[
\alpha-\gamma=\frac{p^2-q^2}{2},
\]
we have
\[
\frac{1}{\alpha}+\frac{p^2-q^2}{2\alpha\gamma}=\frac{1}{\gamma},
\]
and thus
\[
2q^2\alpha\,\Phi'
\ge
\beta^2+2\nu\beta-\frac{\nu\beta^2}{\gamma}.
\]
Subtracting \(2q^2\alpha\) from the right-hand side and simplifying yields
\[
\beta^2+2\nu\beta-\frac{\nu\beta^2}{\gamma}-2q^2\alpha
=
\frac{4\nu^2(4\nu+2p^2+q^2)+q^2(p^2-q^2)^2}{4(2\nu+q^2)}
\ge 0.
\]
Hence
\[
2q^2\alpha\,\Phi'\ge 2q^2\alpha>0,
\]
which proves \(\Phi'(\nu)\ge 1\) on \(I\).

It follows that \(\Phi\) is strictly increasing on \(I\). Since
\[
t=\Phi(\nu(t)),
\]
the function \(t\mapsto \nu(t)\) is nondecreasing. Moreover,
\[
\lambda_3(x,\sqrt t)=x^2+t-\nu(t)=x^2+\Phi(\nu(t))-\nu(t).
\]
Because
\[
\frac{d}{d\nu}\bigl(\Phi(\nu)-\nu\bigr)=\Phi'(\nu)-1\ge 0
\qquad (\nu\in I),
\]
the quantity \(\Phi(\nu)-\nu\) is nondecreasing in \(\nu\), and hence
\[
t\longmapsto \lambda_3(x,\sqrt t)
\]
is nondecreasing on \([0,\infty)\). This proves the proposition.
\end{proof}

\section{The real-slice criterion for \(n=2,3\)}

\begin{theorem}\label{thm:main}
Let \(n\in\{2,3\}\). Then for every \(\eps>0\),
\[
\sigma_\eps(\mathbf A_n)\ \text{is connected}
\quad\Longleftrightarrow\quad
\sigma_\eps(\mathbf A_n)\cap\R\ \text{is connected}.
\]
More precisely, for every \(x\in\R\), the vertical section
\[
V_x^{(n,\eps)}\coloneqq \{y\in\R:\ x+iy\in\sigma_\eps(\mathbf A_n)\}
\]
is either empty or an open interval containing \(0\).
\end{theorem}

\begin{proof}
Fix \(n\in\{2,3\}\), \(\eps>0\), and \(x\in\R\). By Propositions \ref{prop:n2} and
\ref{prop:n3}, the function
\[
|y|\longmapsto \lambda_n(x,y)
\]
is nondecreasing on \([0,\infty)\). Hence, if \(x+iy_0\in \sigma_\eps(\mathbf A_n)\), then
\[
\lambda_n(x,y_0)<\eps^2,
\]
and therefore for every \(y\in\R\) with \(|y|\le |y_0|\),
\[
\lambda_n(x,y)\le \lambda_n(x,y_0)<\eps^2.
\]
Thus
\[
x+iy\in \sigma_\eps(\mathbf A_n).
\]
So each vertical section \(V_x^{(n,\eps)}\) is either empty or an interval containing \(0\).
Since \(\sigma_\eps(\mathbf A_n)\) is open, each such section is in fact an open interval.

Applying Lemma \ref{lem:topological} with \(\Omega=\sigma_\eps(\mathbf A_n)\) gives the desired
connectedness equivalence.
\end{proof}

\begin{remark}
The proofs above establish a stronger fact than connectedness equivalence: for \(n=2\) and
\(n=3\), every vertical slice of the pseudospectrum is centered on the real axis in the sense of
Theorem \ref{thm:main}.
\end{remark}

\end{document}